\chapter{Haptoglobin}
\section*{What Is This Test?}
Haptoglobin is an alpha-2 glycoprotein synthesized by the liver that binds free hemoglobin in the circulation with extremely high affinity (the strongest non-covalent protein-protein interaction known, Kd ~10$^{-15}$ M). The haptoglobin-hemoglobin (Hp-Hb) complex is rapidly cleared from the circulation by CD163 receptors on reticuloendothelial macrophages (primarily in the liver and spleen), preventing the loss of iron (and hemoglobin) through glomerular filtration and protecting the kidneys from hemoglobin-mediated tubular injury \cite{rother2005clinical}. In intravascular hemolysis (e.g., ABO-incompatible transfusion reaction, paroxysmal nocturnal hemoglobinuria [PNH], cold agglutinin hemolytic anemia, mechanical hemolysis from prosthetic heart valves, microangiopathic hemolytic anemia [MAHA], severe G6PD deficiency crisis, or transfusion of incompatible blood), free hemoglobin is released into plasma, depleting haptoglobin (consumed as Hp-Hb complexes are cleared). Serum haptoglobin is the most sensitive and widely used laboratory test for detecting intravascular hemolysis \cite{barcellini2021clinical}. Haptoglobin is decreased in intravascular hemolysis (consumption), decreased in extravascular hemolysis (mild consumption because a small amount of hemoglobin escapes from macrophages into plasma) \cite{packman2019hemolytic}, decreased in impaired hepatic synthesis (cirrhosis, end-stage liver disease), and absent in congenital ahaptoglobinemia. Haptoglobin is an acute phase reactant and may be elevated in active inflammation, infection, malignancy, and trauma; a normal haptoglobin level does NOT exclude hemolysis if the patient has concurrent inflammation.
\section*{Alternative Names}
Haptoglobin; Serum Haptoglobin; Hp; Haptoglobin-Hemoglobin Binding
\section*{Principle}
Haptoglobin is measured by immunoturbidimetry or immunonephelometry using specific anti-haptoglobin antibodies. Anti-human haptoglobin antibodies are incubated with patient serum; the antigen-antibody complex formation increases turbidity proportional to haptoglobin concentration. Results reported in mg/dL \cite{rifai2018tietz}.
\section*{History}
Haptoglobin was discovered in 1938 by Polonovski and Jayle, who identified an alpha-2 plasma globulin that bound hemoglobin and noted that this binding activity disappeared from the serum during hemolysis. In 1955, Oliver Smithies' introduction of starch gel electrophoresis revealed hereditary polymorphism of haptoglobin (types 1-1, 2-1, and 2-2), making it one of the first human protein polymorphisms studied in populations \cite{smithies1955zone}. The binding of hemoglobin to haptoglobin was shown to be extremely tight (dissociation constant on the order of 10$^{-15}$ M) by Hwang and Greer in 1980. The clearance pathway was fully defined in 2001, when Kristiansen and colleagues identified CD163 as the macrophage hemoglobin scavenger receptor \cite{kristiansen2001identification}. Immunoturbidimetric and immunonephelometric assays replaced earlier methods from the 1970s onward, establishing routine clinical measurement of haptoglobin as the most sensitive laboratory marker of intravascular hemolysis.
\section*{Why This Test Is Ordered}
\begin{itemize}
  \item Evaluation of suspected hemolytic anemia, where low haptoglobin supports hemolysis as the cause of anemia
  \item Distinguishing hemolysis from liver disease, bleeding, or ineffective erythropoiesis in a patient with anemia, elevated LDH, and jaundice
  \item Investigation of transfusion reactions and assessment of hemolysis after incompatible blood transfusion
  \item Monitoring of chronic intravascular hemolysis in paroxysmal nocturnal hemoglobinuria, mechanical prosthetic valve hemolysis, and microangiopathic hemolytic anemia
  \item Part of the routine hemolysis panel, ordered with LDH, indirect bilirubin, and reticulocyte count
  \item Assessment of haptoglobin in hemoglobinuria or hemoglobinemia to confirm the intravascular site of hemolysis
\end{itemize}
\section*{Primary vs Secondary Test}
Haptoglobin is a secondary, supporting test in the diagnosis of hemolysis. It is not diagnostic alone and is interpreted within a hemolysis panel (LDH, indirect bilirubin, reticulocyte count, peripheral smear), while excluding the confounding effects of inflammation, liver disease, and recent hemorrhage.
\section*{How This Test Is Performed}
A blood sample is collected by standard venipuncture into a plain serum tube or serum separator tube. Haptoglobin is measured by immunoturbidimetry or immunonephelometry, in which anti-haptoglobin antibodies form immune complexes in proportion to the haptoglobin concentration, with the resulting turbidity or light scattering quantified on an automated analyzer. Results are reported in mg/dL. Turnaround time is typically the same day.
\section*{What Preparation Is Needed}
None; no fasting or special preparation is required. Because haptoglobin falls rapidly once hemolysis begins and can normalize within 3--7 days after a self-limited episode, the sample should be drawn as close as possible to the suspected hemolytic event. It is usually collected together with the rest of the hemolysis panel.
\section*{Contraindications and Precautions}
There are no contraindications beyond those of routine venipuncture. Interpretive cautions: (1) Haptoglobin is an acute phase reactant and rises in inflammation, infection, malignancy, and trauma, so a normal or elevated value does not exclude hemolysis in a patient with concurrent inflammation. (2) Low haptoglobin also occurs with reduced hepatic synthesis (cirrhosis, end-stage liver disease, malnutrition) and in neonates, who have physiologically low levels. (3) Very low haptoglobin can cause spurious increases in unconjugated (indirect) bilirubin measured by some automated methods because of hemoglobin interference. (4) Massive acute hemorrhage can lower haptoglobin transiently through dilution and consumption.
\section*{Risks}
Minimal risk. Slight pain or bruising at the venipuncture site. Rarely: hematoma, infection, or vasovagal syncope.
\section*{What Happens After the Test}
The haptoglobin result is interpreted with LDH, indirect bilirubin, reticulocyte count, and the peripheral smear; the combination of low haptoglobin, elevated LDH, elevated indirect bilirubin, and reticulocytosis confirms hemolytic anemia. The cause of hemolysis is then pursued with directed testing such as the direct antiglobulin test, peripheral smear review, and specific assays for PNH or hemoglobinopathies. In a patient with an acute self-limited episode, haptoglobin may be repeated in 1--2 weeks to document recovery.
\section*{Understanding Results}
\subsection*{Decreased/Absent Haptoglobin (<30 mg/dL)}
Intravascular hemolysis: free hemoglobin >16--20 mg/dL (the binding capacity of haptoglobin) depletes haptoglobin. Causes: ABO-incompatible transfusion, PNH, mechanical valve hemolysis, MAHA (TTP, HUS, DIC), G6PD deficiency crisis, and severe burns. Extravascular hemolysis: Mild to moderate decrease in haptoglobin. Impaired synthesis: Cirrhosis, end-stage liver disease, malnutrition, and neonates (physiologic decrease). Congenital ahaptoglobinemia: Absent haptoglobin by genetic deficiency (absence of the HP gene), more common in African Americans and Asians. Haptoglobin <25 mg/dL is >90\% sensitive and >90\% specific for hemolysis in the absence of liver disease \cite{marchand1980predictive}.
\subsection*{Normal or Elevated Haptoglobin}
No significant hemolysis. Acute phase reactant: infection, inflammation, malignancy, surgery, and trauma. A normal or elevated haptoglobin in a patient with elevated LDH and bilirubin argues against significant hemolysis and suggests an alternative cause (hepatic, biliary, or non-hemolytic).
\section*{Normal Results}
Serum haptoglobin: 30--200 mg/dL. Newborns: Haptoglobin is absent or very low at birth; reaches adult levels by 4--6 months of age. Haptoglobin <25 mg/dL is consistent with hemolysis. The combination of decreased haptoglobin, increased LDH, increased indirect (unconjugated) bilirubin, and increased reticulocyte count confirms hemolytic anemia.
\section*{Follow-up Tests}
Peripheral blood smear (schistocytes in MAHA, spherocytes in autoimmune hemolytic anemia), LDH (markedly elevated in hemolysis and MAHA), indirect bilirubin, reticulocyte count, direct antiglobulin test/DAT (autoimmune hemolytic anemia), hemoglobinemia, hemoglobinuria, and hemosiderinuria (chronic intravascular hemolysis).
\section*{References}
\bibliographystyle{unsrtnat}
\bibliography{references}
\clearpage
